% ============================================
% Response to Reviewer 3
% ============================================
\section*{Response to Reviewer 3}

We thank the reviewer for their rigorous and technically insightful feedback, particularly the geometric analysis of dimensionality effects on similarity metrics. Below we address each point.

% ============================================
\subsection*{R3-1. Clarify L\_1 versus L\_2 normalization in the similarity metric (P1)}

\reviewercomment{With respect to the similarity metrics used by the authors (the dot products $x_A \cdot x_C$ and $x_B \cdot x_C$), I believe that their interpretation is less straightforward than it would appear at first glance. As a first small comment, the authors state that $x$ are ``normalized abundance vectors,'' and Fig. 1B clearly suggests that these are $L_1$ normalized, that is, their elements sum to 1 (i.e., they represent species' relative abundances). But in the Supplementary Methods, the authors clarify that $x$ are $L_2$ normalized, i.e., it is the sum of squares of their elements that is 1. I think making this more clear in the main text would help, since it fundamentally changes what one should expect for the dot products under different hypotheses.}

\responselabel
We thank the reviewer for identifying this ambiguity and apologize for the confusion. To clarify, $\vec{x}$ denotes a community composition vector over ASVs, and the similarity between the coalesced community C and a parent P was computed as $\mathrm{Sim}(C,P) = (\vec{x}_C \cdot \vec{x}_P)/(\|\vec{x}_C\|_2\|\vec{x}_P\|_2)$. Because cosine similarity is invariant to multiplying either vector by a positive scalar, converting each sample from an ASV-abundance vector to a relative-abundance vector before this step gives the same similarity score; the metric depends on the vector direction, not total abundance. We have corrected the Results, Methods, and Supplementary Note~1 wording accordingly. In the experimental analyses, these composition vectors are ASV-abundance vectors derived from the 16S profiles.


Results \S2.1, similarity-score definition: \mschange{``We represent each community by a community composition vector and compute the cosine similarity between the coalesced community C and each parental community.''} The revised equation explicitly divides each dot product by the product of the two Euclidean norms, and the following sentence defines $\vec{x}_A$, $\vec{x}_B$, and $\vec{x}_C$ as community composition vectors.

Methods \S 16S rRNA Sequencing now states: \mschange{``In experimental analyses, the community composition vector was the ASV-abundance vector derived from these 16S profiles.''}

Methods \S Classification of Coalescence Outcomes, at first definition of the similarity metric: \mschange{``Similarity was computed as cosine similarity between community composition vectors (see Supplementary Note~1 for full normalization details).''}

Supplementary Note~1, metric definition paragraph: \mschange{``To characterize coalescence outcomes, we quantified the compositional similarity between the post-coalescence community and each parental community using cosine similarity.''} The note then gives the same equation as the Results and clarifies that sample-wise total-abundance normalization does not change the cosine-similarity score.

% ============================================
\subsection*{R3-2. Dimensionality artifact in similarity metrics (P0)}

\reviewercomment{%
Along these lines, my main concern is that these similarity metrics can behave in somewhat counter-intuitive ways, and some of the results the authors report may be at least partially skewed by simple statistical effects. To make my point as clear as possible, let me propose a simple example. Let's consider two communities A and B, each with N non-overlapping species. Total species' abundances can be represented by vectors $n_A$ and $n_B$:

\[
n_A = ( n_A^{(1)} , n_A^{(2)} , \ldots , n_A^{(N)} , 0 , 0 , \ldots , 0 )
\]
\[
n_B = ( 0 , 0 , \ldots , 0 , n_B^{(1)} , n_B^{(2)} , \ldots , n_B^{(N)} )
\]

Now consider a simple null model (similar to the ``shuffled abundance'' null model proposed by the authors) in which community C, which results from A and B undergoing coalescence, has N species. The abundance vector of C, $n_C$, has N elements that are zero and N that are non-zero. Which elements are set to zero is chosen randomly (though the argument holds in other scenarios, including the ``abundance-weighted random selection'' model formulated by the authors). The non-zero elements of $n_C$ are random samples of the non-zero elements of $n_A$ and $n_B$. This null model aims at representing a ``maximal restructuring'' situation, where species' abundances in community C are completely uncorrelated from their abundances in the parental communities. If we generate $n_A$ and $n_B$ by randomly sampling their non-zero elements uniformly in (0, 1) and we compute the similarity metrics

\[
\mathrm{Sim}(A,C) = n_A \cdot n_C / |n_A| |n_C|
\]
\[
\mathrm{Sim}(B,C) = n_B \cdot n_C / |n_B| |n_C|
\]

For different values of N and different pairs of randomly-generated communities, the ``similarity map'' looks like:

[See attached Reviewer Figure 1]

So despite the model being formulated as a null baseline representing a ``full restructuring'' situation, many cases would be flagged as ``dominance'' or ``mixing.'' The same is true if species' abundances in A and B are sampled differently, for instance if instead of sampling them uniformly ($n \sim U(0,1)$) we used a more uneven distribution (such as $n \sim 10^{U(-3,0)}$) we would find:

[See attached Reviewer Figure 2]

Notably, a somewhat similar effect also occurs under a much simpler null model: if the community C was simply the superposition of communities A and B (i.e., if the two parental communities mixed perfectly without interacting), we would have $n_C = n_A + n_B$. In this case:

[See attached Reviewer Figure 3]

The key takeaway here is that, when using these similarity metrics and the author's ``Dominance-Mixing-Restructuring'' classification, one could expect to observe a decline in the ``Dominance'' fraction as community diversity increases due to a simple statistical effect: positive unit vectors are more likely to be close to orthogonal in lower dimensions than they are in higher dimensions.%
}

\begin{figure}[H]
\centering
\includegraphics[width=0.88\textwidth,height=0.78\textheight,keepaspectratio]{Reviewer3_attachment_figures.pdf}
\caption*{\textbf{Reviewer attachment.}}
\end{figure}

\responselabel
We thank the reviewer for raising this important point. We reproduced the reviewer's toy null models plus the corresponding skewed-abundance additive case (random restructuring with $U(0,1)$, random restructuring with $10^{U(-3,0)}$, simple additive mixing $n_C = n_A + n_B$ with abundances drawn from $U(0,1)$, and simple additive mixing $n_C = n_A + n_B$ with abundances drawn from $10^{U(-3,0)}$) using the exact same L$_2$/cosine analysis used in our manuscript (Response Fig.~R3-2A). We note that one mismatch in the reviewer-attached visualization is that the Restructuring boundary appears to use a retention-magnitude cutoff of $r < 1/2$, whereas our pipeline uses $r \leq 1/\sqrt{2}$ (equivalently, $r^2 \leq 1/2$).

Using the same classification pipeline as in the manuscript, the reproduced results support the reviewer's main concern in a specific regime. The random-restructuring null models become increasingly classified as Restructuring as $N$ grows (uniform case: $43.4\%, 64.9\%, 78.5\%, 88.7\%$ for $N = 2, 4, 6, 8$), and the uniform additive null model concentrates on the symmetric Mixture axis and is classified as Mixture in $80.5\%, 97.4\%, 99.1\%, 99.7\%$ of cases. However, the skewed-abundance random-restructuring null remains high-retention but can fall off the symmetry line and be classified as high Dominance (Dominance: $73.9\%, 50.0\%, 39.7\%, 30.6\%$ for $N = 2, 4, 6, 8$). This confirms that skewed abundance distributions can generate apparent Dominance under simple null models, motivating the event-matched additive-null analysis below.

We have revised the Methods and Supplementary Note~1 to state this boundary explicitly, and Response Fig.~R3-2A shows both the manuscript boundary and the alternative visual boundary at $r = 1/2$.

Methods, ``Classification of Coalescence Outcomes,'' now states that \mschange{``The Restructuring boundary is defined by $r^2 \leq 0.5$, equivalently a retention radius $r \leq 1/\sqrt{2}$ in the similarity map.''} Supplementary Note~1, ``Outcome classification,'' now states that \mschange{``Restructuring occurs when $r^2 \leq 0.5$, equivalently when the retention radius satisfies $r \leq 1/\sqrt{2}$ in the two-parental-community similarity map,''} and \mschange{``the retention threshold $r^2 = 0.5$ (radius $r = 1/\sqrt{2}$, not $r = 1/2$) marks the point at which the best linear combination of the two-parental-community composition vectors accounts for half of the coalesced community's composition in the normalized similarity space.''}

\begin{figure}[H]
\centering
\includegraphics[width=0.95\textwidth]{Fig_R3_2_reviewer_reproduction_L2.pdf}
\caption*{\textbf{Response Fig.~R3-2A.} \textbf{Reviewer toy nulls analyzed with the manuscript classifier.} Rows: \textbf{A}, random restructuring with $U(0,1)$ abundances; \textbf{B}, random restructuring with $10^{U(-3,0)}$ abundances; \textbf{C}, additive mixing with $U(0,1)$ abundances; \textbf{D}, additive mixing with $10^{U(-3,0)}$ abundances. Scatter panels show the reviewer's visual boundary ($r=1/2$, gray) and the manuscript boundaries (retention radius $r=1/\sqrt{2}$ and asymmetry class boundaries, black); bars show manuscript-classifier fractions.}
\end{figure}

\reviewercomment{%
This effect is partially accounted for in the manuscript; specifically, Extended Data Fig.~3 shows that the empirical distribution of Dominance values has significantly higher mean than its counterparts generated by the null models. However, while the authors' claim is correct (``skewed abundance distributions alone do not fully account for the observed asymmetric outcomes''), this does not mean that the empirically quantified Dominance values aren't biased by the statistical effects described above. Many cases flagged as ``Dominance'' may in actuality be compatible with a simple null model, particularly for the communities assembled from the less diverse species pools.%
}

\responselabel
The prior analysis demonstrates that a simple additive null model with skewed abundance can be classified as Dominance because of a simple statistical effect. Specifically, this occurs when the two parental-community abundance vectors have substantially different norms. For a simple additive null model, $n_A+n_B$ is closer to the larger-norm parental community. For non-overlapping parental-community vectors, this effect is direct: $\mathrm{Sim}(A,A+B)$ and $\mathrm{Sim}(B,A+B)$ are proportional to $\|n_A\|_2$ and $\|n_B\|_2$, respectively. We agree that this could particularly matter when abundance is strongly uneven, because vector norms become sensitive to concentration of abundance among a small number of taxa.

However, the norm-imbalance effect is much weaker in our experimental parent pairs than in the skewed toy examples. In our experimental coalescence pairs, the fold differences in Euclidean norms between the two parental raw 16S read-count vectors were $1.85 \pm 0.95$ (mean $\pm$ SD; median $1.58$), with per-medium means ranging from $1.55 \pm 0.51$ to $2.05 \pm 1.18$. These values are far smaller than those in the skewed-abundance simple additive null model, where the fold differences were $18.2 \pm 39.4$ and $7.1 \pm 14.7$ for $N=2$ and $N=4$, respectively (Response Fig.~R3-2B).

To directly test whether this effect could account for the experimental classifications, we constructed a simple additive null model from the actual raw 16S read-count vectors of the same two parental communities in each experimental coalescence pair (Response Fig.~R3-2B). The simple additive null model was classified as Mixture in most events in each nutrient condition: 84/90 in Nutr$-$ (93.3\%), 66/83 in Base (79.5\%), and 69/90 in Nutr$+$ (76.7\%). Thus, the effect identified by the reviewer is present in principle and detectable in some experimental parent pairs, but it is not sufficient to explain the observed Dominance tendency.

\begin{figure}[H]
\centering
\includegraphics[width=0.92\textwidth]{Fig_R3_2_parent_norm_asymmetry.pdf}
\caption*{\textbf{Response Fig.~R3-2B.} \textbf{Parental norm imbalance and simple additive null model classifications.} \textbf{A}, Raw 16S read-count L$_2$ norms for parental-community abundance vectors. \textbf{B}, Paired parental L$_2$ norm fold differences, with skewed-abundance simple additive null models shown for scale; dashed line, $1+\sqrt{2}$ fold difference corresponding to PDI $<0.25$ or $>0.75$ for the simple additive null model. \textbf{C}, Simple additive null model class fractions; labels show Mixture counts.}
\end{figure}

\reviewercomment{%
In my view, it would be necessary to compare each coalescence community with a corresponding null model constructed specifically for each of the pairs of parental communities, such as the simple additive model $n_{C,\mathrm{null}} = n_A + n_B$. Comparing the experimental coalesced community with such a null model should be done on a case-by-case basis, not just by comparing distributions of Dominance indexes as in Extended Data Fig.~3.%
}

\responselabel
We thank the reviewer for this suggestion. We compared each observed Base-medium coalescence community with the corresponding simple additive null model, constructed from the same two parental communities using raw 16S read-count vectors. As described above, this null model was usually classified as Mixture, although it produced Dominance in a subset of events. Critically, the observed Base-medium outcomes were shifted further toward direction-independent parental asymmetry, $y = |2\mathrm{PDI}-1|$, than their matched additive-null outcomes (77/83 events; one-sided paired Wilcoxon test, $p = 5.9 \times 10^{-14}$). Thus, the event-matched comparison shows that the simple additive model can contribute to apparent asymmetry in some cases but does not by itself explain the observed Dominance tendency. We replaced Extended Data Fig.~3 with this event-matched additive-null analysis and described the associated norm-imbalance effect in Supplementary Note~1.

\begin{figure}[H]
\centering
\includegraphics[width=0.98\textwidth]{Fig_R3_2_additive_null_comparison.pdf}
\caption*{\textbf{Response Fig.~R3-2C.} \textbf{Base-medium raw-count observations versus the event-matched simple additive null model.} \textbf{A}, Raw-count similarity map: orange open points, simple additive null model outcomes ($n_{C,\mathrm{null}}=n_A+n_B$); blue filled points, observed coalesced communities; red arrows, five example matched null-to-observation pairs. \textbf{B}, Paired shifts in direction-independent parental asymmetry, $y = |2\mathrm{PDI}-1|$, from additive null to observation for all 83 Base events; gray lines connect matched pairs and the dashed line marks the $y=0.5$ asymmetry boundary. \textbf{C}, Class transitions from additive-null outcomes (columns) to observations (rows).}
\end{figure}

% ============================================
\subsection*{R3-3. Interaction strength, diversity, and Dominance frequency (P0)}

\reviewercomment{In section 2.3 of the Results, the authors examine how varying the mean interaction (competition) strength between species in a gLV model affects coalescence outcomes. The main conclusion here is that increasing competition strength shifts the distribution of coalescence outcomes from all ``Mixing'' to mostly ``Dominance'' and ``Restructuring.'' However, it is not clear to what extent this phenomenon can be explained by simple statistical effects like the ones described above. Increasing the mean interaction strength (mu) can be expected to decrease the number of species in the communities, and, as mentioned above, reducing N generically causes an increase in the frequency of ``Dominance'' outcomes. It seems important to test whether the increase of ``Dominance'' outcomes in the model (and in the experiments) as mu grows is substantially different from the expected increase in the null models as N decreases.}

\responselabel
As shown in the R3-2 analyses, in our experiment the richness and norm-imbalance effect identified by the reviewer can occur, but it is not sufficient to explain the observed Dominance increase. For the experimental data, the simple additive null model shows that paired parental raw 16S read-count vector norm differences were relatively small and that the simple additive null model was usually classified as Mixture rather than Dominance (Response Fig.~R3-2B). Thus, this statistical effect contributes to the null expectation in some events but does not account for the experimental Dominance pattern.

In the gLV simulations, increasing the interaction-strength parameter reduces final richness and creates greater unevenness in parental-community norms, so Dominance outcomes in the simple additive null model become more frequent because of this richness- and norm-imbalance effect (Response Fig.~R3-3A). However, this effect is insufficient to explain Dominance in the full simulated coalescence outcomes. At the three representative interaction-strength parameter values used in the main figure, the simple additive null model gives $0.0\%$, $5.7\%$, and $16.2\%$ Dominance at $\mu = 0.3$, $0.6$, and $0.8$, respectively, whereas the full gLV coalescence simulations give $20.4\%$, $61.0\%$, and $69.5\%$ Dominance. Thus, richness- and norm-imbalance effects contribute to the null expectation at high $\mu$, but they do not explain the simulated Dominance transition.

We have added this analysis to Supplementary Note~5, in the paragraph on abundance-skew and richness-dependent geometric effects, and show the full analysis in Supplementary Fig.~38.

\begin{figure}[H]
\centering
\includegraphics[width=0.90\textwidth]{Fig_R3_3_sim_parent_norm_asymmetry.pdf}
\caption*{\textbf{Response Fig.~R3-3A.} \textbf{Simulation parental-community composition-vector norm imbalance and simple additive null model classifications across interaction strength.} \textbf{A}, Euclidean norms of assembled parental-community composition vectors across $\mu$; line shows the median and shading shows the 25th--75th percentile range. \textbf{B}, Fold difference in Euclidean norm between the two parental communities for each simulated coalescence event; line shows the median and shading shows the 25th--75th percentile range, and dotted line shows the 95th percentile. The dashed horizontal line marks the fold difference $1+\sqrt{2}$ corresponding to PDI $<0.25$ or $>0.75$ for the simple additive null model. \textbf{C}, Classification fractions for the simple additive null model computed from the same paired parental-community composition vectors at each $\mu$.}
\end{figure}

\reviewercomment{It would also be useful to see how richness (in the final communities, not in the inocula) changes with mu in the model, as well as across the different media (Base, Nutr-, Nutr+) used in the experiment. From Supplementary Figs. 10-12, it seems like communities derived from natural samples do exhibit substantial diversity differences in the Base and Nutr+ media with respect to the Nutr- medium, but I could not find an equivalent representation for the first experiment.}

\responselabel
We agree that showing the final-community richness directly is useful. In the simulations, final richness decreases with $\mu$ for both post-assembly sub-communities and coalesced communities (Response Fig.~R3-3B, panel A). In the experiment, final coalesced-community richness also differs across media, with mean richness of 13.0, 9.9, and 8.5 ASVs in Nutr$-$, Base, and Nutr$+$, respectively (Response Fig.~R3-3B, panel B). Thus, richness differences are real and are now shown directly. However, as demonstrated by the prior analyses, the associated statistical/geometric effect does not explain the Dominance patterns: it does not account for the experimental event-matched additive-null result (Response Fig.~R3-2B,C) or for the simulation trend (Response Fig.~R3-3A).

\begin{figure}[H]
\centering
\includegraphics[width=0.92\textwidth]{Fig_R3_3_richness_summary.pdf}
\caption*{\textbf{Response Fig.~R3-3B.} \textbf{Final richness across interaction strength and nutrient conditions.} \textbf{A}, Mean final richness of simulated post-assembly sub-communities and coalesced communities across $\mu$; error bars show s.e.m. \textbf{B}, Final coalesced-community richness (ASVs above 0.1\% threshold) across nutrient conditions and initial pool sizes in the synthetic-community experiment.}
\end{figure}

% ============================================
\subsection*{R3-4. gLV model excludes facilitation}

\reviewercomment{With respect to the model, the gLV framework is restricted to competitive interactions ($\alpha_{ij} > 0$). While this is reasonable from a modeling perspective because it prevents situations of ``unchecked growth,'' it also limits the range of ecological scenarios that the model can represent. Previous works have shown that facilitative interactions (particularly in the form of metabolic cross-feeding) are an important determinant of correlated species' fates during coalescence.}

\responselabel
We thank the reviewer for raising this important and valid limitation. We treat competition as the major source of interaction in this experimental system and in our modeling, while recognizing that facilitation can be important in other microbial coalescence settings. This approximation is supported by the 12-isolate pairwise coculture assay, where focal ASV yields were usually lower in coculture than in monoculture: 79/84/93\,\% of species-in-pair observations fell below the monoculture expectation in Nutr$-$/Base/Nutr$+$ (Response Fig.~R3-4a).

We agree that the competition-only gLV model cannot capture facilitative interactions such as metabolic cross-feeding. To consider the broader ecological context, we ran the same coalescence analysis for two model extensions: (1) facilitative tails that allow positive ecological interactions (negative $\alpha_{ij}$ under our sign convention), and (2) reciprocal sign/correlation structure between coefficients $\alpha_{ij}$ and $\alpha_{ji}$. The simulation definitions and full parameter sweeps are now described in Supplementary Note~3 and Supplementary Figs.~39--40.

In the facilitative-tail extension, we preserved the mean interaction coefficient and added density-dependent self-limitation to prevent runaway growth (Response Fig.~R3-4b). Dominance still rises with $\mu$ at every facilitative-tail strength tested. At $\mu=0.30$, increasing the facilitative tail shifts outcomes away from Mixture toward both Dominance and Restructuring (Dom/Mix/Res $=18/73/9\,\%$ at $f=0$ and $43/33/24\,\%$ at $f=0.8$). At $\mu=0.80$, Dominance remains high across the facilitative tail (70--76\%). Thus, facilitative interactions can change the balance between Mixture and Restructuring in some regimes, but the core claim, the monotonic rise of Dominance with stronger interaction ($\mu$), remains largely valid in this model extension.

We also varied reciprocal pair-coupling structure, using antisymmetric exploitation for $p<0$ and symmetric competition for $p>0$ (Response Figs.~R3-4c,d). This analysis shows that reciprocal coefficient correlation does not affect our core claim: Dominance still increases with stronger interaction ($\mu$) when reciprocal coefficients are anticorrelated or positively correlated.

We clarified this scope in the Supplementary Methods: \mschange{``Because we restrict $\alpha_{ij} \geq 0$, the model captures only competitive (growth-inhibiting) interactions; facilitative interactions such as cross-feeding are not represented.''} The Discussion now states the corresponding limitation: \mschange{``Additionally, both our experimental system and theoretical model focus on a competition-dominated regime; systems with substantial mutualism or facilitation may exhibit qualitatively different dynamics.''}

\begin{figure}[H]
\centering
\includegraphics[width=\textwidth]{R3_4_experiment.pdf}
\caption*{\textbf{Response Fig.~R3-4a.} \textbf{Pairwise coculture data show widespread growth-inhibiting effects, especially in Base and Nutr$+$.}
\textbf{(A--C)} Per-medium scatter of each focal ASV's coculture CFU ($C_i$) versus its monoculture CFU ($M_i$) from the 12-isolate pairwise invasion assay. Each point is one ASV in one unordered pair, averaged across the two invasion directions when both were available; the dashed line marks $C_i=M_i$. Points below the diagonal indicate direct suppression of that ASV in coculture relative to monoculture.
\textbf{(D)} Bar summary of the fraction of species-in-pair observations below the monoculture expectation: 79/84/93\,\% for Nutr$-$/Base/Nutr$+$. These fractions are descriptive summaries; no p-values are shown for this panel.}
\label{fig:R3-4a}
\end{figure}

\begin{figure}[H]
\centering
\includegraphics[width=\textwidth]{R3_4_mixed_sign_higher_order.pdf}
\caption*{\textbf{Response Fig.~R3-4b.} \textbf{Allowing a facilitative tail can increase Restructuring but does not remove the rise of Dominance with $\mu$.}
Off-diagonal coefficients are drawn iid from $\alpha_{ij}\sim U[-f\mu,(2+f)\mu]$, preserving $E[\alpha_{ij}]=\mu$ while increasing the expected facilitative fraction $P(\alpha_{ij}<0)$ from 0 to 22.2\%. Top insets show the analytic sampling density of $\alpha_{ij}/\mu$ (red: facilitative support; gray: nonnegative support), drawn from the formula rather than from simulated matrices. Dynamics include a stabilizing density-dependent self-limitation term, $\dot n_i=n_i(1-(A n)_i-0.1n_i^2)$. Each panel fixes one facilitative-tail strength $f$ and shows Dominance / Mixture / Restructuring fractions at $\mu \in \{0.30,0.60,0.80\}$ (200 pools $\times$ 6 coalescence events per cell). Trends across the simulated parameter grid are descriptive; no p-values are shown for this robustness simulation.}
\label{fig:R3-4b}
\end{figure}

\begin{figure}[H]
\centering
\includegraphics[width=\textwidth]{R3_4_simulation.pdf}
\caption*{\textbf{Response Fig.~R3-4c.} \textbf{Dominance still increases with $\mu$ across reciprocal pair-coupling structures.}
Each panel fixes one $p$ and shows Dominance / Mixture / Restructuring fractions at $\mu \in \{0.30, 0.60, 0.80\}$, ordered from antisymmetric exploitation ($p<0$) through the i.i.d.\ competitive baseline ($p=0$) to symmetric competition ($p>0$). Top insets show the analytic directed-coefficient distribution implied by each $p$ value for $\alpha_{ij}/\mu$ (red: negative coefficients; gray: nonnegative coefficients), drawn from the sampler definition rather than from Monte Carlo samples. Data: 200 pools $\times$ 6 coalescence events per cell. Trends across $\mu$ and $p$ are descriptive; no p-values are shown for this robustness simulation.}
\label{fig:R3-4c}
\end{figure}

\begin{figure}[H]
\centering
\includegraphics[width=\textwidth]{R3_4_pair_coupling_fine.pdf}
\caption*{\textbf{Response Fig.~R3-4d.} \textbf{The finer pair-coupling sweep preserves the $\mu$-dependent Dominance trend.}
Same pair-coupling definition as Response Fig.~R3-4c, but with $p$ swept from $-1$ to $+1$ in increments of 0.2. Top strips show representative analytic directed-coefficient distributions for $p=-1,0,+1$; the sweep itself uses the same sampler for all $p$ values. Each panel fixes $\mu$ and shows stacked Dominance / Mixture / Restructuring fractions across $p$; the black line marks the Dominance fraction. Trends across the simulated $p$ sweep are descriptive; no p-values are shown for this robustness simulation.}
\label{fig:R3-4d}
\end{figure}

\reviewercomment{Would two species which engage in a mutualism be considered ``weak interactors'' or the opposite?}

\responselabel
Within our framework, interaction strength refers to coefficient magnitude and distributional spread (the May-style random-Lotka--Volterra axis; May 1972; Grilli 2017), not interaction sign. Thus, two species in a strong mutualism would be strong interactors, with coefficients of opposite sign from strong competitors. We note that our empirical invasion-resistance assay, however, measures growth-inhibiting effects and would not score mutualistic facilitation as strong competitive exclusion.

To test whether this sign distinction affected the model result, we added weak bidirectional mutualistic pairs to the gLV simulations. Specifically, 0--40\% of unordered species pairs were made bidirectionally mutualistic ($\alpha_{ij},\alpha_{ji}\sim U[-0.2\mu,0]$), while the remaining interactions were left competitive. The $\mu$-dependent Dominance trend persisted: Dominance remained 58.5--70.2\% at $\mu=0.60$ and 69.4--77.5\% at $\mu=0.80$ (Response Fig.~R3-4e; Supplementary Fig.~41). Thus, adding weak mutualistic pairs does not remove the Dominance trend.

This magnitude/spread framing is further supported by the mean-vs-variance gLV sweep $\alpha_{ij}\sim U[m-h,m+h]$ with $m,h \in \{0,0.2,0.4,0.6,0.8\}$: zero coefficient half-width ($h=0$) gave only Mixture outcomes even at the highest mean ($m=0.8$), and across the 25 cells Dominance correlated more strongly with the half-width $h$ than with the mean $m$ (descriptive Pearson $r=0.78$ vs.\ $0.53$; Response Fig.~R3-4f; Supplementary Fig.~42). Mechanistically, heterogeneous coefficients are what let assembly filter species into mutually compatible subgroups whose members face stronger competition from outsiders, the ``cohesion without cooperation'' route to Dominance noted in the Discussion (Tikhonov 2016; May 1972; Grilli 2017). We clarified this mean-vs-variance decomposition in Supplementary Methods (``Lotka--Volterra Simulations'') and Supplementary Note~3, also in response to Reviewer 2's related request for biological interpretation of $\mu$ and the coefficient distributions (R2-6).

\begin{figure}[H]
\centering
\includegraphics[width=\textwidth]{R3_4_mutualistic_pair_fraction.pdf}
\caption*{\textbf{Response Fig.~R3-4e.} \textbf{Weak bidirectional mutualistic pairs do not remove the $\mu$-dependent Dominance trend.}
For a controlled fraction $q \in \{0,0.10,0.20,0.30,0.40\}$ of unordered species pairs, both reciprocal coefficients are sampled from $\alpha_{ij},\alpha_{ji}\sim U[-0.2\mu,0]$, corresponding to weak positive ecological interactions under our sign convention. All remaining off-diagonal coefficients are sampled from the baseline competitive distribution $U[0,2\mu]$. Dynamics include the same stabilizing density-dependent self-limitation term used in Response Fig.~R3-4b, $\dot n_i=n_i(1-(A n)_i-0.1n_i^2)$. Each panel fixes one mutualistic-pair fraction $q$ and shows Dominance / Mixture / Restructuring fractions at $\mu \in \{0.30,0.60,0.80\}$ (200 pools $\times$ 6 coalescence events per cell). Trends across the simulated parameter grid are descriptive; no p-values are shown for this robustness simulation.}
\label{fig:R3-4e}
\end{figure}

\begin{figure}[H]
\centering
\includegraphics[width=0.92\textwidth]{R3_4_mean_variance_grid.pdf}
\caption*{\textbf{Response Fig.~R3-4f.} \textbf{Coefficient variance, not coefficient mean alone, drives Dominance in the mixed-sign gLV sweep.}
Off-diagonal coefficients are sampled from $\alpha_{ij}\sim U[m-h,m+h]$, with mean coefficient $m$ and standard deviation $h/\sqrt{3}$ varied independently. Negative $\alpha_{ij}$ values, corresponding to positive ecological interactions under our sign convention, occur above the dashed $h=m$ line. Dynamics include the same stabilizing density-dependent self-limitation term used in Response Fig.~R3-4b, $\dot n_i=n_i(1-(A n)_i-0.1n_i^2)$. Heatmaps show Dominance, Mixture, Restructuring, and same-parent concordance $|\phi|$ across the 5 $\times$ 5 grid (200 pools $\times$ 6 coalescence events per cell). Trends are descriptive; no p-values are shown for this robustness simulation.}
\label{fig:R3-4f}
\end{figure}

\reviewercomment{This seems particularly relevant given that ``Restructuring'' appears to be much more common in natural sample-derived communities than in the initial bottom-up-assembled communities. Given these discrepancies, and the fact that natural sample-derived communities are still model systems (even if they may be closer to natural microbiomes than communities assembled from defined species pools), I would suggest the claims in section 2.6 of the Results as well as in the Discussion are toned down (e.g., lines 292-294).}

\responselabel
We have changed and toned down our conclusion for natural sample-derived communities in Results \S2.6: \mschange{``These communities showed higher Restructuring fractions, possibly reflecting greater taxonomic diversity and more complex interaction networks. These results suggest that the nutrient-dependent increase in Dominance observed in synthetic consortia is also present in taxonomically richer natural sample-derived communities.''} We also removed the detailed natural-community generality caveat from the Discussion and consolidated it in Supplementary Note~6, which emphasizes that these are laboratory-stabilized communities rather than unfiltered natural ecosystems: \mschange{``Thus, pre-selection could contribute to the convergent coalescence patterns observed between synthetic and natural sample-derived communities, and the natural-community results should be interpreted as evidence that the nutrient-dependent increase in Dominance occurs in taxonomically richer laboratory-stabilized communities, not as evidence for unrestricted generality across unfiltered natural ecosystems.''}

% ============================================
\subsection*{R3-5. ``Interaction strength'' versus ``competition strength''}

\reviewercomment{One last comment along these lines, which is admittedly mostly semantic: the term ``interaction strength'' may be somewhat confusing since, in the context of the gLV model, it refers strictly to competitive interactions. It may be worth considering using different terminology (``competition strength''?)}

\responselabel
We thank the reviewer for raising this important concern about the semantics around the term ``interaction strength''; after revisiting the terminology, we retain this term because it is consistent with the random-interaction and microbial gLV literature that motivates our framework, including May's classic random community model (May 1972) and recent microbial gLV work by Hu et al. (2022, 2025), where interaction strength is used prevalently. The R3-4 analyses support this scope by showing that the empirical coculture data contain frequent growth-inhibiting effects, that the qualitative Dominance trend persists in facilitative-tail and pair-coupling model extensions, and that coefficient variance is required when coefficient mean and variance are separated. We therefore kept the established term while clarifying that, in the baseline gLV model, $\mu$ sets the range of sampled competitive interaction coefficients and controls both their mean and variance. Methods now states: \mschange{``We fixed growth rates to 1 and self-interaction coefficients $\alpha_{ii} = 1$, and drew off-diagonal competition coefficients from a uniform distribution $\mathbb{U}(0, 2\mu)$, so $\mu$ sets the range of sampled competitive interaction coefficients and thereby controls both their mean and variance within this phenomenological model.''} We also clarified in the revised Discussion, quoted above in R3-4, that the model focuses on a competitive-interaction regime.
